% Supplementary D — Journal §IV (Method)
% Full mathematical derivations for each axis F-1 through F-12.

\documentclass[11pt,a4paper]{article}
% Shared preamble for all supplementary chapter documents.
% Used by each supplementary/conference/suppl_*.tex and
% supplementary/journal/suppl_*.tex via % Shared preamble for all supplementary chapter documents.
% Used by each supplementary/conference/suppl_*.tex and
% supplementary/journal/suppl_*.tex via % Shared preamble for all supplementary chapter documents.
% Used by each supplementary/conference/suppl_*.tex and
% supplementary/journal/suppl_*.tex via \input{../preamble.tex}.

\usepackage[utf8]{inputenc}
\usepackage[T1]{fontenc}
\usepackage[a4paper,margin=2.2cm]{geometry}
\usepackage{amsmath,amssymb,amsfonts,amsthm}
\usepackage{mathtools}
\usepackage{graphicx}
\usepackage{booktabs}
\usepackage{array}
\usepackage{multirow}
\usepackage{longtable}
\usepackage{algorithm}
\usepackage{algpseudocode}
\usepackage{xcolor}
\usepackage{url}
\usepackage{cite}
\usepackage[hidelinks]{hyperref}

\graphicspath{{../../figures/}}

\renewcommand{\thesection}{\Alph{section}.\arabic{section}}
\renewcommand{\thesubsection}{\thesection.\arabic{subsection}}

\theoremstyle{plain}
\newtheorem{theorem}{Theorem}[section]
\newtheorem{lemma}[theorem]{Lemma}
\newtheorem{proposition}[theorem]{Proposition}
\theoremstyle{definition}
\newtheorem{definition}[theorem]{Definition}
\theoremstyle{remark}
\newtheorem{remark}[theorem]{Remark}

\setcounter{secnumdepth}{3}
\setcounter{tocdepth}{2}

\newcommand{\R}{\mathbb{R}}
\newcommand{\E}{\mathbb{E}}
\newcommand{\KL}[2]{\mathrm{KL}\!\left(#1 \,\middle\|\, #2\right)}
\newcommand{\ari}{\mathrm{ARI}}
\newcommand{\nmi}{\mathrm{NMI}}
\newcommand{\softmax}{\mathrm{softmax}}
.

\usepackage[utf8]{inputenc}
\usepackage[T1]{fontenc}
\usepackage[a4paper,margin=2.2cm]{geometry}
\usepackage{amsmath,amssymb,amsfonts,amsthm}
\usepackage{mathtools}
\usepackage{graphicx}
\usepackage{booktabs}
\usepackage{array}
\usepackage{multirow}
\usepackage{longtable}
\usepackage{algorithm}
\usepackage{algpseudocode}
\usepackage{xcolor}
\usepackage{url}
\usepackage{cite}
\usepackage[hidelinks]{hyperref}

\graphicspath{{../../figures/}}

\renewcommand{\thesection}{\Alph{section}.\arabic{section}}
\renewcommand{\thesubsection}{\thesection.\arabic{subsection}}

\theoremstyle{plain}
\newtheorem{theorem}{Theorem}[section]
\newtheorem{lemma}[theorem]{Lemma}
\newtheorem{proposition}[theorem]{Proposition}
\theoremstyle{definition}
\newtheorem{definition}[theorem]{Definition}
\theoremstyle{remark}
\newtheorem{remark}[theorem]{Remark}

\setcounter{secnumdepth}{3}
\setcounter{tocdepth}{2}

\newcommand{\R}{\mathbb{R}}
\newcommand{\E}{\mathbb{E}}
\newcommand{\KL}[2]{\mathrm{KL}\!\left(#1 \,\middle\|\, #2\right)}
\newcommand{\ari}{\mathrm{ARI}}
\newcommand{\nmi}{\mathrm{NMI}}
\newcommand{\softmax}{\mathrm{softmax}}
.

\usepackage[utf8]{inputenc}
\usepackage[T1]{fontenc}
\usepackage[a4paper,margin=2.2cm]{geometry}
\usepackage{amsmath,amssymb,amsfonts,amsthm}
\usepackage{mathtools}
\usepackage{graphicx}
\usepackage{booktabs}
\usepackage{array}
\usepackage{multirow}
\usepackage{longtable}
\usepackage{algorithm}
\usepackage{algpseudocode}
\usepackage{xcolor}
\usepackage{url}
\usepackage{cite}
\usepackage[hidelinks]{hyperref}

\graphicspath{{../../figures/}}

\renewcommand{\thesection}{\Alph{section}.\arabic{section}}
\renewcommand{\thesubsection}{\thesection.\arabic{subsection}}

\theoremstyle{plain}
\newtheorem{theorem}{Theorem}[section]
\newtheorem{lemma}[theorem]{Lemma}
\newtheorem{proposition}[theorem]{Proposition}
\theoremstyle{definition}
\newtheorem{definition}[theorem]{Definition}
\theoremstyle{remark}
\newtheorem{remark}[theorem]{Remark}

\setcounter{secnumdepth}{3}
\setcounter{tocdepth}{2}

\newcommand{\R}{\mathbb{R}}
\newcommand{\E}{\mathbb{E}}
\newcommand{\KL}[2]{\mathrm{KL}\!\left(#1 \,\middle\|\, #2\right)}
\newcommand{\ari}{\mathrm{ARI}}
\newcommand{\nmi}{\mathrm{NMI}}
\newcommand{\softmax}{\mathrm{softmax}}


\title{Journal Supplementary D --- Full mathematical derivations\\
{\large Companion to: \emph{Beyond Accuracy} (Section IV)}}
\author{Felipe Santib\'a\~nez-Leal}
\date{2026}

\begin{document}
\maketitle

\section{Scope}
This supplement provides the complete mathematical content for the
twelve-axis framework: generative models, variational inference,
hierarchical Bayesian likelihood, coherence formulae, ARI and
Hungarian definitions, and the algorithmic pseudocode for each axis.

\section{Canonical generative model}

\subsection{LDA on band-frequency documents}
For $K$ topics, vocabulary $V = B$ band indices, $D$ documents and
hyperparameters $\alpha, \eta > 0$:
\begin{align}
\phi_k    & \sim \mathrm{Dir}(\eta\mathbf{1}_V), \quad k=1,\ldots,K, \\
\theta_d  & \sim \mathrm{Dir}(\alpha\mathbf{1}_K), \quad d=1,\ldots,D, \\
z_{d,n}   & \sim \mathrm{Cat}(\theta_d), \quad n=1,\ldots,N_d, \\
w_{d,n}   & \sim \mathrm{Cat}(\phi_{z_{d,n}}).
\end{align}

\subsection{Online VB lower bound}
The mean-field variational family is
$q(\theta_d, z_d, \phi) = q(\theta_d \mid \gamma_d)
\big(\prod_n q(z_{d,n} \mid \lambda_{d,n})\big) \prod_k q(\phi_k \mid
\beta_k)$. The per-document ELBO is
\begin{align}
\mathcal{L}_d
&= \sum_n \sum_k \lambda_{d,n,k}
\big[
   \psi(\gamma_{d,k}) - \psi(\textstyle\sum_{k'} \gamma_{d,k'})
   \nonumber\\
&\quad + \psi(\beta_{k, w_{d,n}}) - \psi(\textstyle\sum_v \beta_{k,v})
   - \log \lambda_{d,n,k}
\big]
\nonumber\\
&\quad + (\alpha-1) \sum_k\!\big[
   \psi(\gamma_{d,k}) - \psi(\textstyle\sum_{k'}\gamma_{d,k'})\big]
\nonumber\\
&\quad + \log \Gamma(\textstyle\sum_{k'} \gamma_{d,k'})
       - \sum_k \log \Gamma(\gamma_{d,k}) + \mathrm{const}.
\end{align}
Updates are Eq.~(8) and Eq.~(9) of the conference supplement B.
Online VB applies the natural-gradient ascent
$\beta \leftarrow (1 - \rho_t) \beta + \rho_t \tilde{\beta}_t$
with $\rho_t = (\tau_0 + t)^{-\kappa}$, $\tau_0 = 10$, $\kappa = 0.7$.

\subsection{Posterior summaries}
$\phi_{k,b} = \beta_{k,b} / \sum_v \beta_{k,v}$;
$\theta_{d,k} = \gamma_{d,k} / \sum_{k'} \gamma_{d,k'}$;
$z_d = \arg\max_k \theta_{d,k}$.

\section{Axis F-1: hierarchical Bayesian model}

\subsection{Likelihood}
For $M$ methods, $S$ scenes, $F$ folds, observed fold-OA scores
$y_{m,s,f}$:
\begin{align}
y_{m,s,f}      & = \mu_m + \delta_s + \rho_f + \epsilon_{m,s,f}, \\
\mu_m          & \sim \mathcal{N}(0, 1), \\
\delta_s       & \sim \mathcal{N}(0, 0.5^2), \\
\rho_f         & \sim \mathcal{N}(0, 0.2^2), \\
\epsilon_{m,s,f} & \sim \mathcal{N}(0, \sigma^2), \\
\sigma         & \sim \mathrm{HalfNormal}(0.5).
\end{align}

\subsection{Inference}
NUTS with 1000 draws + 1000 tune, 2 chains, $\hat{R} < 1.01$ on
all parameters in all reported runs. Posterior mean and 94\% HDI of
$\mu_m$ are reported; pairwise $P[\mu_A > \mu_B]$ are estimated from
the joint posterior draws as
\begin{equation}
P[\mu_A > \mu_B] \approx \frac{1}{T} \sum_{t=1}^{T}
\mathbf{1}[\mu_A^{(t)} > \mu_B^{(t)}],
\end{equation}
with $T = 2000$ draws across two chains.

\section{Axis F-2: coherence measures}

\subsection{$c_v$ coherence}
For a topic $t$ with top-$N$ words $W_t = (w_1, \ldots, w_N)$, the
NPMI between words $w_i, w_j$ on a sliding window of size $s$ is
\begin{equation}
\mathrm{NPMI}(w_i, w_j) = \frac{\log
   \frac{P(w_i, w_j) + \epsilon}{P(w_i) P(w_j)}
}{- \log (P(w_i, w_j) + \epsilon)}.
\end{equation}
The topic context vector is
$\mathbf{v}_t = \sum_{i=1}^{N} \mathbf{u}_i$ with
$u_{i,j} = \mathrm{NPMI}(w_i, w_j)$. The $c_v$ score is
\begin{equation}
c_v(t) = \frac{1}{N} \sum_{i=1}^{N}
\cos(\mathbf{v}_t, \mathbf{u}_i),
\end{equation}
the average cosine similarity between the topic context vector and
each word's individual context vector. We use $N = 15$ and the
default Röder sliding window of size $s = 110$.

\subsection{NPMI and U-Mass}
$c_{\text{NPMI}}(t) = \binom{N}{2}^{-1} \sum_{i<j} \mathrm{NPMI}(w_i,
w_j)$.
$\mathrm{U\!-\!Mass}(t) = \binom{N}{2}^{-1} \sum_{i<j} \log
\frac{N(w_i, w_j) + 1}{N(w_j)}$.

\section{Axis F-3: seed stability}
For each (scene, method): refit five LDA / ProdLDA / ETM under seeds
$\{0, \ldots, 4\}$. For each refit compute KMeans clustering of
$\theta$ into $K$ clusters; report ARI($z_{\text{KMeans}}$, label),
$c_v$. Aggregate across seeds as mean $\pm$ std.

\section{Axis F-4: capacity sensitivity}
Refit LDA at $K \in \{4, 6, 8, 10, 12, 16, 20, 24\}$ with the
canonical hyperparameters, $N = 5$ seeds per $K$. Per-$K$ aggregates:
\begin{itemize}
\item perplexity$_K^{\text{test}} = \exp(-\frac{1}{|\mathcal{D}_{\text{test}}|}
\sum_{d \in \mathcal{D}_{\text{test}}} \log p(\mathbf{w}_d \mid \phi_K, \alpha))$,
\item topic-diversity$_K = \frac{|\bigcup_k W_k^{(K)}|}{N \cdot K}$
with $W_k^{(K)}$ the top-$N$ words per topic at capacity $K$,
\item matched-cosine$_K = $ mean over (seed pair, topic) of
$\cos(\phi^{(s_1)}_k, \phi^{(s_2)}_{\sigma^*(k)})$.
\end{itemize}

\section{Axis F-5: paired ARI under masks}
Detailed in conference supplement B; restated here for cross-reference.
$\ari_{\text{paired}}(m) = \ari(\{z_d\}_\mathcal{S},
\{z^{(m)}_d\}_\mathcal{S})$.

\section{Axis F-6: cross-method ARI}
For each scene compute the partition set $\mathcal{P} = \{\text{label},
\text{topic-dom}, \text{Felz}, \text{SLIC-500}, \text{SLIC-2000},
\text{p-15}, \text{p-7}, \text{px}\}$. The reported matrix is
\begin{equation}
M_{ij} = \ari(\mathcal{P}_i, \mathcal{P}_j),
\end{equation}
restricted to pixels with a valid (non-sentinel) value in both
partitions. NMI matrix is computed analogously.

\section{Axis F-7: topic--label coupling}
For each (canonical topic $t$, label $\ell$),
$P(\ell \mid t) = N_{t,\ell} / \sum_{\ell'} N_{t,\ell'}$ with
$N_{t,\ell}$ the count of pixels with dominant topic $t$ and ground
truth $\ell$. Per-topic Shannon entropy:
$H(L \mid t) = - \sum_\ell P(\ell \mid t) \log P(\ell \mid t)$.
Per-topic KL divergence from the marginal:
$\KL{P(L \mid t)}{P(L)} = \sum_\ell P(\ell \mid t)
\log [P(\ell \mid t) / P(\ell)]$.

\section{Axis F-8: per-topic Hungarian alignment}
Confusion matrix
$C \in \mathbb{Z}^{K \times K^{(m)}}_{\ge 0}$,
$C_{k, j} = |\{d \in \mathcal{S} : z_d = k \wedge z^{(m)}_d = j\}|$.
Hungarian: $\sigma^* = \arg\max_\sigma \sum_k C_{k, \sigma(k)}$,
solved by \texttt{scipy.optimize.linear\_sum\_assignment} on $-C$.
Swap rate:
$\mathrm{swap} = |\mathcal{S}|^{-1} \sum_d \mathbf{1}[z^{(m)}_d \ne
\sigma^*(z_d)]$.

\section{Axis F-9: HIDSAG preprocessing stability}
For each HIDSAG subset, $R = 3$ preprocessing recipes $\{\text{cont-rem},
\text{deriv}, \text{smooth}\}$. For each recipe pair $(r_a, r_b)$
compute $\ari(\text{dom-top}_{r_a}, \text{dom-top}_{r_b})$ on the
documents present in both recipes. Per-subset aggregate is the mean
over the $\binom{3}{2} = 3$ pairs.

\section{Axis F-10: cross-scene transfer}
Given scene $s$ with canonical basis $\phi^{(s)}$, project the
documents of scene $t$ through one E-step
\begin{equation}
\gamma_{d, k}^{(t \leftarrow s)} = \alpha + \sum_n
\lambda_{d, n, k}^{(t \leftarrow s)},
\end{equation}
with $\lambda$ computed from $\phi^{(s)}$ frozen. The transfer
report is
$\ari(\text{label}_t, \arg\max_k \theta^{(t \leftarrow s)})$.

\section{Axis F-11: rate--distortion}
Reconstruct
$\hat{\mathbf{x}}_d = \sum_k \theta_{d,k} \phi_k^\top$ in the
band-frequency-count space, then de-quantise to a continuous
reflectance via the inverse of the canonical tokenisation
$\hat{x}_{d,b} = c_{d,b}^{\text{recon}} / Q$. Distortion is
$\mathrm{MSE} = D^{-1} \sum_d \|\hat{\mathbf{x}}_d
- \mathbf{x}_d\|^2$ at rate $\log_2 K$ bits.

\section{Axis F-12: external baseline}
For each labelled scene, collect three reference OA values from
prior published work. Report
$\Delta_s = \text{OA}_s^{\text{topic-routed-soft}}
- \max_{r \in \text{refs}(s)} \text{OA}_s^{(r)}$
with the per-paper split caveat.

\section{Inverse Hungarian (used in F-8 swap-rate computation)}
The post-permutation confusion matrix $\tilde{C}$ used for the
journal Figure 3(b) panel is
\begin{equation}
\tilde{C}_{k, j} = C_{k, \sigma^{*-1}(j)},
\end{equation}
which moves the matched diagonal to the main diagonal. The right
panel of journal Figure 3 visualises $\tilde{C}$ directly.

\bibliographystyle{IEEEtran}
\bibliography{../../bibliography/refs}

\end{document}
